% !TEX root = ../main.tex
\section{Introduction}
\label{sec:introduction}

Stem diameters in forest stands are commonly characterised by fitting statistical distributions to plot measurements \citep{bailey1973quantifying, hyink1983generalized}. Because permanent sample plot programmes are expensive to maintain, analysts often rely on legacy inventories whose tally rules prioritised field efficiency over statistical convenience.

Merchantability thresholds remove small stems, biological limits cap large diameters, and crews aggregate counts into fixed-width classes, leaving doubly truncated grouped stand tables.

Complete-form distributions defined on $(0, \infty)$ perform poorly on such data because tail bias persists and residual sums fail to vanish. Analysts therefore evaluate truncated versions of the target distribution \citep{zutter1986characterizing}, yet closed forms are not always published, forcing teams to reproduce earlier derivations \citep{mittal1989estimating, hannon1999estimation, hegde1989estimation, ahmad2001moments, wingo1989left} or numerically integrate $F(x; \boldsymbol{P}) = \int_0^x f(t)\,\mathrm{d}t$.

We propose a two-stage method that retains the \emph{complete} form while achieving truncated-quality fits. The method is benchmarked against one-stage complete (1sc) and one-stage truncated (1st) fits using permanent sample plot data from Qu\'ebec, Canada. Across Weibull and gamma families the two-stage complete (2sc) estimator matches specialised truncated fits while preserving implementation simplicity.
